\section{Background}

The concept of recursive spawning in multi-agent systems draws from several converging research traditions: lifecycle management of autonomous agents, accountability and provenance in distributed agentic deployments, the distinction between fixed and mutable delegation chains, hierarchical task decomposition, and the architectural substrate that makes immutable parent-pointer enforcement tractable. This section reviews each area in turn, establishing the conceptual ground from which the BX3 Framework's recursive spawning model emerges.

%---------------------------------------------------------------
\subsection{Agent Lifecycle Management in Multi-Agent Systems}

Modern multi-agent systems require explicit lifecycle management: the ability to reason about when an agent is created, what authority it holds, how it operates, and when it is terminated or transferred. The IETF draft on Agent Identity, Trust and Lifecycle Protocol (AITLP) \cite{draft-larsson-aitlp} identifies five lifecycle stages---creation, onboarding, operation, revocation, and legacy transfer---and emphasizes that each transition must preserve the integrity of the agent's mandate and its connection to the authorizing principal. AITLP introduces the concept of an Agent Manifest: a formal declaration of identity, scope, and commitment that travels with an agent throughout its lifecycle, enabling downstream systems to verify legitimacy without relying on centralized registries.

Complementing AITLP, the human-anchored identity architecture draft \cite{draft-beyer-agent-identity-architecture} argues that accountability in multi-agent systems must trace back to a durable human identity root. Rather than treating agents as independently authorized entities, this architecture requires that every agent action be attributable to a human principal through an explicit delegation chain. The architecture distinguishes between \emph{explicit delegation semantics}---where the scope, constraints, and conditions of authority are formally expressed---and implicit or ad-hoc authorization patterns that dominate most current agent deployments. This distinction is foundational: without explicit delegation, the provenance of an agent's actions cannot be reliably audited.

Practical accountability models for multi-agent organizations have been formalized in the JaCaMo-based framework by~\citeA{multiagent-accountability}, who define an accountability lifecycle that governs how obligations are created, monitored, reported, and adjudicated. Their model introduces the concept of a \texttt{Delay(G)} event---an organizational signal indicating that an obligation related to goal \textit{G} has gone unfulfilled---which propagates up the authority hierarchy until it reaches a responsible agent. This model demonstrates that accountability is not a static property but an active, propagating process that must be architecturally supported rather than assumed.

%---------------------------------------------------------------
\subsection{Accountability and Provenance in Distributed Systems}

As agentic workflows span multiple execution environments---edge devices, cloud services, and HPC基础设施---provenance tracking becomes essential for transparency, reproducibility, and regulatory compliance. PROV-AGENT \cite{prov-agent} extends the W3C PROV standard with agent-centric metadata, capturing prompts, responses, decisions, and tool executions as first-class provenance objects. PROV-AGENT integrates with the Model Context Protocol (MCP) to link individual agent decisions to the full end-to-end workflow, enabling fine-grained lineage tracking and fault localization. In evaluation across diverse infrastructures, PROV-AGENT demonstrated practical utility for near real-time provenance capture and query-based reliability analysis, offering a concrete path toward auditable multi-agent deployments.

The Human Delegation Provenance (HDP) protocol \cite{draft-hdp-agentic-delegation} addresses a specific gap in existing token-based schemes: the need to cryptographically record and verify human authorization across multi-hop delegation chains. HDP tokens bind a human authorization event to an AI session, then record each agent's handling of that token as a signed hop in an append-only chain. Verification requires only the issuer's Ed25519 public key and the current session identifier---no network lookups or third-party trust anchors. This design enables both runtime integrity checking (detecting prompt injection by verifying the chain of authorization) and post-hoc audits for regulatory compliance.

Contextual Agent Authorization Mesh (CAAM) \cite{draft-barney-caam} further extends authorization semantics by introducing a Session Context Object (SCO) that binds delegated authority to a specific purpose, session, and trust chain. The SCO encodes constraints that persist across agent-to-agent interactions, ensuring that authority delegated in one context cannot be repurposed in another. CAAM's post-discovery authorization handshake provides a runtime mechanism for enforcing these constraints after agent discovery but before tool execution.

For enterprise-grade deployments, the trace-based assurance framework of \citeA{trace-based-assurance} captures agentic executions as Message-Action Traces (MAT), with explicit step-level and trace-level contracts yielding machine-checkable verdicts. This framework provides quantitative metrics---task success, termination reliability, contract compliance, factuality indicators, containment rate, and governance outcome distributions---that enable rigorous comparison of multi-agent orchestration designs under stochastic conditions.

%---------------------------------------------------------------
\subsection{Fixed vs. Mutable Delegation Chains}

A critical and under-explored distinction in multi-agent delegation is whether the chain of delegation is \emph{fixed} or \emph{mutable} after creation. Mutable delegation chains allow a parent agent to retroactively modify the scope or target of a previously issued delegation, which simplifies certain dynamic reallocation scenarios but introduces significant accountability risks: a downstream action may become unauthorized after the fact, yet have already executed based on a now-invalid delegation token.

The BX3 Framework adopts a \emph{fixed delegation chain} model. When a parent node births a child loop, the delegation token is cryptographically sealed at creation time. The child operates under immutable authorization for the duration of its scope. This design choice trades dynamic reallocation flexibility for a stronger accountability guarantee: any action taken by the child is attributable to the exact delegation token that existed at the moment of authorization, with no possibility of retroactive invalidation by a parent acting on updated information.

This fixed-chain model aligns with the append-only provenance philosophy espoused by HDP \cite{draft-hdp-agentic-delegation}, where delegation hops are recorded rather than modified. It also contrasts with OAuth 2.0 Token Exchange and JSON Web Tokens, which support in-band token refresh and replacement without necessarily preserving the original delegation context \cite{draft-hdp-agentic-delegation}. The BX3 model's immutability guarantee is architecturally enforced at the Purpose Layer: the parent cannot modify the child delegation without spawning a new child with a fresh token.

The architectural implication is that delegation chains in BX3-compliant systems are \emph{growth-only}: new children can be created, but existing delegation tokens remain intact. This property enables reliable auditability, supports the Bailout Protocol's upward-escalation mechanism, and ensures that the provenance chain from any action back to its originating human principal is reconstructable at any future point.

%---------------------------------------------------------------
\subsection{Hierarchical Task Decomposition in AI Agents}

Hierarchical task decomposition is the process by which a complex goal is recursively broken into subgoals, each of which may be further decomposed until atomic actions are reached. This approach is foundational to classical planning formalisms, notably Hierarchical Task Networks (HTNs) \cite{chathtm}, where decomposition is performed by matching task symbols against pre-defined methods and operators in a domain knowledge base. The key property of HTN planners is \emph{soundness}: any plan they produce is guaranteed to achieve the specified task, even when some decompositions are approximated by language models \cite{chathtm}.

Recent work has extended hierarchical decomposition to LLM-driven agents operating in long-horizon, stochastic environments. STRUCTUREDAGENT \cite{structuredagent} introduces dynamic AND/OR trees for multi-level planning and recursive decomposition of complex tasks, coupled with a structured memory module that maintains candidate solutions and tracks constraint satisfaction. ReAcTree \cite{reactree} organizes subgoals in dynamically constructed agent trees, where each node is an LLM agent capable of reasoning, acting, and further expanding the tree. ReAcTree demonstrates that hierarchical agent trees substantially outperform flat reasoning baselines on long-horizon benchmarks---approximately doubling goal success rates in some settings.

Task-Decoupled Planning (TDP) \cite{task-decoupled-planning} replaces monolithic, entangled reasoning with a directed acyclic graph (DAG) of sub-goals, enabling modular planning and execution where each sub-task operates within a scoped context. TDP reduces token usage by up to 82\% compared to entangled approaches while improving robustness to execution errors, as local corrections do not disrupt the overall workflow. The Supervisor-Planner-Executor architecture in TDP formalizes the roles that BX3 assigns to Purpose, Bounds Engine, and Fact Layer respectively.

ReCAP \cite{recap} addresses the context drift problem in long-horizon LLM planning by combining plan-ahead decomposition with structured re-injection: at each depth level, the agent generates a full subtask plan, executes the first item, then re-injects the remaining parent plan into the LLM context as a sliding window. This mechanism maintains coherence between high-level intent and low-level actions across multiple recursion levels while keeping context requirements linear in task depth.

Deep Agent \cite{deep-agent} introduces the Hierarchical Task DAG (HTDAG), which dynamically decomposes high-level goals into interdependent sub-tasks across multiple layers. HTDAG supports autonomous API and tool creation, a recursive two-stage planner-executor that continuously refines plans, and real-time graph modification in response to new information. This aligns with BX3's recursive spawning model, where child loops maintain local goal structures while remaining subordinate to the parent's accountability.

%---------------------------------------------------------------
\subsection{The BX3 Framework as Substrate for Immutable Parent Chains}

The BX3 Framework, introduced by \citeA{bx3-whitepaper-v2}, defines three immutable functional layers that collectively enforce the accountability and provenance guarantees required for recursive spawning. The \textbf{Purpose Layer} provides intent, judgment, and accountability; it is the layer at which goals are set, trade-offs are weighed, and final accountability is assigned. Critically, the Purpose Layer does not execute actions---it defines what ought to happen and is responsible for what does happen. Any actor occupying the Purpose Layer must satisfy the layer's functional requirements, regardless of whether that actor is a human, an AI system, or a mechanical process \cite{bx3-whitepaper-v2}.

The \textbf{Bounds Engine} provides interpretation, bounded reasoning, and constrained execution. It analyzes, recommends, and decides within its bounded domain, but any action that crosses into the physical world is mediated by the Fact Layer. The \textbf{Fact Layer} provides deterministic enforcement, hard physical constraint, and forensic auditability. It is the only layer that can produce non-reversible physical effects, and its defining characteristic is 100\% deterministic behavior: given identical inputs, it must produce identical outputs \cite{bx3-whitepaper-v2}.

Recursive spawning in the BX3 Framework operates as follows: a parent node births a child loop that carries a hard-coded pointer to the parent's Purpose Layer. Each child loop operates within a containerized \emph{Worksheet} that preserves local state and local goal structures while remaining subordinate to the parent's accountability. The delegation token is fixed at creation time and cannot be modified by the parent retroactively. This design is the architectural realization of the fixed delegation chain principle described above.

The BX3 framework's Loop Isolation property ensures that the Bounds Engine and Fact Layer never share a functional plane, meaning a single AI agent cannot simultaneously propose and execute an action \cite{bx3-whitepaper-v2}. This architectural constraint eliminates the most common failure mode in mutable-delegation systems: an agent that modifies its own authorization scope mid-execution.

The BX3 Accountability Invariant states that every action in the Fact Layer is traceable to a Purpose Layer actor. The Escalation Invariant states that any state the Bounds Engine cannot resolve within its current authority triggers the Bailout Protocol, reaching a human anchor without passing through intermediate machine actors. The Determinist Invariant guarantees that the Fact Layer produces identical outputs for identical inputs, enabling formal verification and forensic audit \cite{bx3-whitepaper-v2}. Together, these invariants provide the foundational guarantees upon which the recursive spawning model's immutable parent chains rest.

The Training Wheels Protocol \cite{training-wheels} provides the operational mechanism for enforcing the accountability invariant during the recursive spawning process. Under Mode 1 (HITL Active), every outbound action from a spawned child is queued for human review before execution, ensuring that the fixed delegation chain does not produce unauthorized physical effects before the system has demonstrated trustworthy behavior. This progressive, earned-autonomy model is what makes the BX3 recursive spawning substrate practically defensible in production deployments.

The separation of concerns in BX3 maps cleanly onto the architectural requirements identified by the trace-based assurance framework \cite{trace-based-assurance}: per-agent capability limits, action mediation at the language-to-action boundary, and containment guarantees under realistic faults. The BX3 Fact Layer's deterministic constraint enforcement is structurally equivalent to the trace-level contracts that yield machine-checkable compliance verdicts in that framework.

The governance translation work of \citeA{koch2026}---translating governance norms into enforceable runtime controls---further supports the BX3 approach of encoding accountability guarantees in architectural invariants rather than relying on policy documentation. The type-checked compliance model for agentic financial systems \cite{rashie2026} demonstrates a complementary enforcement path using formal verification in the Lean 4 theorem prover, suggesting that the BX3 immutable parent chain model can be formally verified in safety-critical domains.

The convergence of these research threads---lifecycle management protocols, cryptographic provenance chains, hierarchical decomposition methods, and functional-layer architectures---establishes both the need for and the feasibility of the BX3 recursive spawning model. The following sections formalize the model, characterize its properties, and demonstrate its application in the Agentic runtime.